\documentclass[12pt,a4paper]{amsart}
\usepackage{amsmath,amssymb,amsthm}
\usepackage{mathtools}
\usepackage[utf8]{inputenc}
\usepackage[T1]{fontenc}
\usepackage{hyperref}
\usepackage{enumitem,booktabs,array}
\usepackage{geometry}
\geometry{margin=2.5cm}

\newtheorem{theorem}{Theorem}[section]
\newtheorem{lemma}[theorem]{Lemma}
\newtheorem{corollary}[theorem]{Corollary}
\newtheorem{proposition}[theorem]{Proposition}
\newtheorem{conjecture}[theorem]{Conjecture}
\theoremstyle{definition}
\newtheorem{definition}[theorem]{Definition}
\newtheorem{example}[theorem]{Example}
\newtheorem{remark}[theorem]{Remark}

\author{Michael Fuhrmann}
\address{Specialist Mathematics Project, Build 141}
\email{specialist-maths@localhost}
\date{\today}

\begin{document}
\title{Special Functions in Mathematics and Physics:\\
Gamma, Bessel, Hypergeometric, and Related Functions}

\begin{abstract}
Special functions arise as canonical solutions to the most fundamental differential
equations of mathematical physics and analysis. This paper provides a systematic
treatment of the principal special functions: the Gamma and Beta functions, Bessel
functions of the first and second kind, the Gauss hypergeometric function \({}_2F_1\),
Legendre polynomials and associated Legendre functions, Hermite and Laguerre polynomials
with their quantum-mechanical applications, elliptic functions of Weierstra\ss{} and
Jacobi, and the Riemann zeta function together with Dirichlet \(L\)-functions.
For each family we establish the defining differential equation or series, prove
fundamental identities such as the reflection formula \(\Gamma(z)\Gamma(1-z)=\pi/\sin(\pi z)\),
the Beta--Gamma relation, and the functional equation of \(\zeta(s)\), and discuss
connections among these function families.  Numerical verifications drawn from the
accompanying Python module \texttt{special\_functions.py} corroborate all theoretical
results.
\end{abstract}

\maketitle

\tableofcontents

%% ============================================================
\section{Introduction}
%% ============================================================

Special functions are those functions that arise so frequently in the solution of canonical
problems in analysis, geometry, and mathematical physics that they have been given
individual names and extensive tables of properties.  Their systematic study constitutes
one of the most beautiful chapters in classical analysis.

The central objects of this paper are solutions of second-order linear ordinary
differential equations with polynomial coefficients.  By Fuchs's theorem, such equations
possess power-series solutions near regular singular points, and the resulting functions
display rich algebraic and analytic structure: recurrence relations, generating functions,
orthogonality on an appropriate interval, asymptotic expansions, and connections to
number theory.

\begin{example}[Canonical differential equations]
The following second-order ODEs each define a fundamental family of special functions:
\begin{align}
  x^2 y'' + x y' + (x^2 - \nu^2)y &= 0 \qquad \text{(Bessel equation)},\label{eq:bessel_ode}\\
  (1-x^2)y'' - 2xy' + n(n+1)y &= 0 \qquad \text{(Legendre equation)},\label{eq:legendre_ode}\\
  z(1-z)w'' + [c-(a+b+1)z]w' - ab\,w &= 0 \qquad \text{(hypergeometric equation)},\label{eq:hyp_ode}\\
  y'' &= x\,y \qquad \text{(Airy equation)}.\label{eq:airy_ode}
\end{align}
\end{example}

The remainder of this paper is organised as follows.
Section~\ref{sec:gamma} treats the Gamma function, including its reflection and
duplication formulas.  Section~\ref{sec:beta} covers the Beta function and its integral
representation.  Section~\ref{sec:bessel} develops Bessel functions.
Section~\ref{sec:hypergeometric} discusses the hypergeometric function \({}_2F_1\).
Section~\ref{sec:legendre} presents Legendre polynomials and spherical harmonics.
Section~\ref{sec:hermite_laguerre} covers Hermite and Laguerre polynomials and their role
in quantum mechanics.  Section~\ref{sec:elliptic} treats elliptic functions.
Section~\ref{sec:zeta} discusses the Riemann zeta function and Dirichlet \(L\)-functions.
Section~\ref{sec:applications} surveys physical applications.

%% ============================================================
\section{The Gamma Function}
\label{sec:gamma}
%% ============================================================

\subsection{Definition and Basic Properties}

\begin{definition}[Gamma function via integral]
For \(\Re(z) > 0\) define
\[
  \Gamma(z) \;=\; \int_0^\infty t^{z-1} e^{-t}\,dt.
\]
\end{definition}

Integration by parts immediately yields the fundamental functional equation.

\begin{theorem}[Functional equation]\label{thm:gamma_functional}
For all \(z\) with \(\Re(z) > 0\),
\[
  \Gamma(z+1) = z\,\Gamma(z).
\]
In particular, \(\Gamma(n+1) = n!\) for every \(n \in \mathbb{N}_0\).
\end{theorem}

\begin{proof}
Integrate by parts:
\[
  \Gamma(z+1)
  = \int_0^\infty t^z e^{-t}\,dt
  = \bigl[-t^z e^{-t}\bigr]_0^\infty + z\int_0^\infty t^{z-1}e^{-t}\,dt
  = 0 + z\,\Gamma(z). \qquad\qquad\square
\]
\end{proof}

The functional equation together with \(\Gamma(1) = 1\) gives the meromorphic continuation
of \(\Gamma\) to all of \(\mathbb{C}\): \(\Gamma\) has simple poles at
\(z = 0, -1, -2, \ldots\) and no zeros.

\subsection{Weierstra\ss{} Product and Gauss Formula}

The Euler--Mascheroni constant \(\gamma = -\Gamma'(1) \approx 0.5772\) appears in the
Weierstra\ss{} product representation
\[
  \frac{1}{\Gamma(z)}
  = z\,e^{\gamma z}\prod_{n=1}^\infty \left(1+\frac{z}{n}\right)e^{-z/n}.
\]
Equivalently, the Gauss limit formula reads
\[
  \Gamma(z) = \lim_{n\to\infty} \frac{n!\,n^z}{z(z+1)\cdots(z+n)}.
\]

\subsection{Reflection Formula}

\begin{theorem}[Euler reflection formula]\label{thm:gamma_reflection}
For all \(z \notin \mathbb{Z}\),
\[
  \Gamma(z)\,\Gamma(1-z) = \frac{\pi}{\sin(\pi z)}.
\]
\end{theorem}

\begin{proof}
Using the Weierstra\ss{} product for \(1/\Gamma(z)\) and \(1/\Gamma(1-z)\),
\[
  \frac{1}{\Gamma(z)\,\Gamma(-z)}
  = -z\prod_{n=1}^\infty\left(1-\frac{z^2}{n^2}\right).
\]
The infinite product for \(\sin(\pi z)/(\pi z)\) gives
\[
  \frac{\sin(\pi z)}{\pi z} = \prod_{n=1}^\infty\left(1-\frac{z^2}{n^2}\right),
\]
whence \(\Gamma(z)\Gamma(-z) = -\pi/(z\sin(\pi z))\).
Replacing \(z\mapsto 1-z\) and using \(\Gamma(1-z) = -z\,\Gamma(-z)\) yields the result.
\end{proof}

\begin{corollary}
\(\Gamma\!\left(\tfrac{1}{2}\right) = \sqrt{\pi}\).
\end{corollary}
\begin{proof}
Set \(z = 1/2\) in Theorem~\ref{thm:gamma_reflection}:
\(\Gamma(1/2)^2 = \pi/\sin(\pi/2) = \pi\).
\end{proof}

\subsection{Duplication Formula}

\begin{theorem}[Legendre duplication formula]\label{thm:gamma_duplication}
For all \(z\) with \(\Re(z) > 0\),
\[
  \Gamma(z)\,\Gamma\!\left(z+\tfrac{1}{2}\right)
  = \frac{\sqrt{\pi}}{2^{2z-1}}\,\Gamma(2z).
\]
\end{theorem}

\begin{proof}
Consider the Beta integral \(B(z,z) = \Gamma(z)^2/\Gamma(2z)\).  The substitution
\(t = (1+u)/2\) transforms this into
\[
  B(z,z) = 2^{1-2z}\int_{-1}^1 (1-u^2)^{z-1}\,du = 2^{1-2z} B\!\left(\tfrac{1}{2},z\right)
  = 2^{1-2z}\frac{\sqrt{\pi}\,\Gamma(z)}{\Gamma(z+1/2)},
\]
giving the formula upon rearrangement.
\end{proof}

\subsection{Stirling's Approximation}

For large \(|z|\) with \(|\arg z| < \pi\),
\[
  \ln\Gamma(z) = \left(z-\tfrac{1}{2}\right)\ln z - z + \tfrac{1}{2}\ln(2\pi)
  + \sum_{k=1}^N \frac{B_{2k}}{2k(2k-1)z^{2k-1}} + O\!\left(z^{-2N-1}\right),
\]
where \(B_{2k}\) are Bernoulli numbers.  In particular,
\[
  n! \;\sim\; \sqrt{2\pi n}\,\left(\frac{n}{e}\right)^n \quad (n\to\infty).
\]

%% ============================================================
\section{The Beta Function}
\label{sec:beta}
%% ============================================================

\begin{definition}[Beta function]
For \(\Re(a) > 0\), \(\Re(b) > 0\),
\[
  B(a,b) = \int_0^1 t^{a-1}(1-t)^{b-1}\,dt.
\]
\end{definition}

\begin{theorem}[Beta--Gamma relation]\label{thm:beta_gamma}
\[
  B(a,b) = \frac{\Gamma(a)\,\Gamma(b)}{\Gamma(a+b)}.
\]
\end{theorem}

\begin{proof}
Write \(\Gamma(a)\Gamma(b)\) as a double integral:
\[
  \Gamma(a)\Gamma(b)
  = \int_0^\infty s^{a-1}e^{-s}\,ds\int_0^\infty t^{b-1}e^{-t}\,dt.
\]
Substituting \(s = r\,u\), \(t = r(1-u)\) with \(r \in (0,\infty)\), \(u \in (0,1)\)
and Jacobian \(r\) gives
\[
  = \int_0^\infty r^{a+b-1}e^{-r}\,dr \cdot \int_0^1 u^{a-1}(1-u)^{b-1}\,du
  = \Gamma(a+b)\,B(a,b). \qquad\square
\]
\end{proof}

\begin{proposition}[Symmetry]
\(B(a,b) = B(b,a)\).
\end{proposition}

\begin{proposition}[Trigonometric form]
\[
  B(a,b) = 2\int_0^{\pi/2}\sin^{2a-1}\theta\,\cos^{2b-1}\theta\,d\theta.
\]
\end{proposition}

\begin{remark}
The Beta function appears naturally in the computation of volumes of balls in
\(\mathbb{R}^n\), in the theory of the Dirichlet distribution, and as normalisation
constants for orthogonal polynomials.
\end{remark}

%% ============================================================
\section{Bessel Functions}
\label{sec:bessel}
%% ============================================================

\subsection{The Bessel Differential Equation}

\begin{definition}[Bessel equation of order \(\nu\)]
\begin{equation}\label{eq:bessel}
  x^2 y'' + x y' + (x^2 - \nu^2)y = 0, \qquad x > 0,\; \nu \in \mathbb{C}.
\end{equation}
\end{definition}

Equation~\eqref{eq:bessel} has a regular singular point at \(x = 0\) with indicial roots
\(\pm\nu\).  The Frobenius method yields two linearly independent solutions denoted
\(J_\nu(x)\) and \(Y_\nu(x)\) (for non-integer \(\nu\); for integer \(\nu\) a logarithmic
term appears in \(Y_\nu\)).

\subsection{Bessel Functions of the First Kind}

\begin{definition}[Bessel function \(J_\nu\)]
\[
  J_\nu(x) = \sum_{m=0}^\infty \frac{(-1)^m}{m!\,\Gamma(m+\nu+1)}
  \left(\frac{x}{2}\right)^{2m+\nu}.
\]
\end{definition}

For integer \(n \in \mathbb{Z}\) the integral representation holds:
\[
  J_n(x) = \frac{1}{\pi}\int_0^\pi \cos(n\tau - x\sin\tau)\,d\tau.
\]

\subsection{Bessel Functions of the Second Kind (Neumann Functions)}

\begin{definition}[Neumann function \(Y_\nu\)]
For non-integer \(\nu\),
\[
  Y_\nu(x) = \frac{J_\nu(x)\cos(\nu\pi) - J_{-\nu}(x)}{\sin(\nu\pi)}.
\]
For integer \(n\), \(Y_n(x) = \lim_{\nu\to n} Y_\nu(x)\), which involves \(\ln x\).
\end{definition}

\(Y_\nu(x)\) is singular at \(x = 0\); in physical problems involving the origin one
retains only \(J_\nu\).

\subsection{Modified Bessel Functions}

The substitution \(x \mapsto ix\) in~\eqref{eq:bessel} yields the modified Bessel equation
\(x^2 y'' + xy' - (x^2+\nu^2)y = 0\) whose solutions are the modified Bessel functions
\[
  I_\nu(x) = i^{-\nu}J_\nu(ix), \qquad
  K_\nu(x) = \frac{\pi}{2}\,\frac{I_{-\nu}(x)-I_\nu(x)}{\sin(\nu\pi)}.
\]
\(K_\nu(x)\) decays exponentially: \(K_\nu(x) \sim \sqrt{\pi/(2x)}\,e^{-x}\) as
\(x \to +\infty\).

\subsection{Recurrence Relations}

\begin{theorem}[Three-term recurrences]\label{thm:bessel_recur}
For \(x \neq 0\):
\begin{align}
  J_{\nu-1}(x) + J_{\nu+1}(x) &= \frac{2\nu}{x}\,J_\nu(x),\label{eq:bessel_rec1}\\
  J_{\nu-1}(x) - J_{\nu+1}(x) &= 2J_\nu'(x).\label{eq:bessel_rec2}
\end{align}
\end{theorem}

\subsection{Wronskian}

\begin{theorem}[Wronskian of \(J_\nu\) and \(Y_\nu\)]
\[
  W\bigl[J_\nu(x), Y_\nu(x)\bigr]
  = J_\nu(x)Y_\nu'(x) - J_\nu'(x)Y_\nu(x)
  = \frac{2}{\pi x}.
\]
\end{theorem}

This non-vanishing Wronskian confirms that \(J_\nu\) and \(Y_\nu\) form a fundamental
system for \(x > 0\).

\subsection{Zeros and Asymptotic Behaviour}

All zeros of \(J_\nu\) for \(\nu > -1\) are real and simple.  Asymptotically,
\[
  J_\nu(x) \sim \sqrt{\frac{2}{\pi x}}\cos\!\left(x - \frac{\nu\pi}{2} - \frac{\pi}{4}\right)
  \quad (x \to +\infty).
\]

%% ============================================================
\section{Hypergeometric Functions}
\label{sec:hypergeometric}
%% ============================================================

\subsection{The Gauss Hypergeometric Series}

\begin{definition}[Pochhammer symbol]
\[
  (a)_0 = 1, \qquad (a)_n = a(a+1)\cdots(a+n-1) = \frac{\Gamma(a+n)}{\Gamma(a)}.
\]
\end{definition}

\begin{definition}[Gauss hypergeometric function]\label{def:2F1}
For \(c \notin \{0,-1,-2,\ldots\}\) and \(|z| < 1\),
\[
  {}_2F_1(a,b;c;z)
  = \sum_{n=0}^\infty \frac{(a)_n(b)_n}{(c)_n\,n!}\,z^n.
\]
\end{definition}

The series converges absolutely for \(|z| < 1\) and can be analytically continued to
\(\mathbb{C}\setminus[1,\infty)\).

\subsection{Gauss Hypergeometric Differential Equation}

\begin{theorem}
\({}_2F_1(a,b;c;z)\) satisfies the equation
\[
  z(1-z)w'' + \bigl[c-(a+b+1)z\bigr]w' - ab\,w = 0.
\]
This equation has regular singular points at \(z = 0, 1, \infty\).
\end{theorem}

\subsection{Transformation Formulas}

\begin{theorem}[Pfaff transformation]
\[
  {}_2F_1(a,b;c;z) = (1-z)^{-a}\,{}_2F_1\!\left(a,c-b;c;\frac{z}{z-1}\right).
\]
\end{theorem}

\begin{theorem}[Euler transformation]
\[
  {}_2F_1(a,b;c;z) = (1-z)^{c-a-b}\,{}_2F_1(c-a,c-b;c;z).
\]
\end{theorem}

\begin{theorem}[Kummer's 24 solutions]
Kummer found 24 solutions of the hypergeometric equation obtained by applying the six
symmetries of the triple \((0,1,\infty)\) and the two local solutions at each point.
\end{theorem}

\subsection{Special Cases}

\begin{proposition}[Important specialisations]
\begin{align*}
  {}_2F_1(1,1;1;z) &= \frac{1}{1-z} \quad (|z|<1),\\
  {}_2F_1\!\left(\tfrac{1}{2},\tfrac{1}{2};1;k^2\right) &= \frac{2}{\pi}\,K(k),\\
  P_n(x) &= {}_2F_1\!\left(-n,n+1;1;\tfrac{1-x}{2}\right),\\
  \ln(1+z) &= z\,{}_2F_1(1,1;2;-z).
\end{align*}
\end{proposition}

\subsection{Confluent Hypergeometric (Kummer) Functions}

Taking \(b \to \infty\) with \(z \mapsto z/b\) in Definition~\ref{def:2F1} produces the
confluent hypergeometric function (Kummer function):
\[
  M(a;c;z) = {}_1F_1(a;c;z)
  = \sum_{n=0}^\infty \frac{(a)_n}{(c)_n\,n!}\,z^n,
\]
satisfying the Kummer equation \(z w'' + (c-z)w' - aw = 0\).  Notably,
\(M(a;a;z) = e^z\) and
\(\operatorname{erf}(z) = (2z/\sqrt{\pi})\,M(1/2;3/2;-z^2)\).

%% ============================================================
\section{Legendre Polynomials}
\label{sec:legendre}
%% ============================================================

\subsection{Definition via Rodrigues' Formula}

\begin{theorem}[Rodrigues formula]\label{thm:rodrigues}
\[
  P_n(x) = \frac{1}{2^n n!}\,\frac{d^n}{dx^n}(x^2-1)^n.
\]
\end{theorem}

The first few polynomials are:
\[
  P_0(x)=1,\quad P_1(x)=x,\quad P_2(x)=\tfrac{1}{2}(3x^2-1),\quad
  P_3(x)=\tfrac{1}{2}(5x^3-3x).
\]

\subsection{Orthogonality}

\begin{theorem}[Orthogonality of Legendre polynomials]\label{thm:legendre_ortho}
\[
  \int_{-1}^1 P_n(x)\,P_m(x)\,dx = \frac{2}{2n+1}\,\delta_{nm}.
\]
\end{theorem}

\begin{proof}
Both \(P_n\) and \(P_m\) satisfy Legendre's equation.  Multiply the equation for \(P_n\)
by \(P_m\), subtract the analogous product for \(P_m\) times \(P_n\), and integrate over
\([-1,1]\).  The boundary terms vanish because \((1-x^2)\) is zero at \(\pm 1\), yielding
\([n(n+1)-m(m+1)]\int_{-1}^1 P_n P_m\,dx = 0\).  For \(n \neq m\) this forces the
integral to zero.  The normalisation \(2/(2n+1)\) follows from \(\|P_n\|^2\) computed via
Rodrigues' formula and integration by parts \(n\) times.
\end{proof}

\subsection{Generating Function}

\begin{theorem}[Generating function]
For \(|t| < 1\) and \(x \in [-1,1]\),
\[
  \frac{1}{\sqrt{1-2xt+t^2}} = \sum_{n=0}^\infty P_n(x)\,t^n.
\]
\end{theorem}

\subsection{Associated Legendre Functions and Spherical Harmonics}

\begin{definition}[Associated Legendre function]
For \(0 \leq m \leq n\),
\[
  P_n^m(x) = (-1)^m(1-x^2)^{m/2}\frac{d^m}{dx^m}P_n(x).
\]
\end{definition}

\begin{definition}[Spherical harmonics]
\[
  Y_n^m(\theta,\varphi)
  = \sqrt{\frac{(2n+1)}{4\pi}\frac{(n-|m|)!}{(n+|m|)!}}\,
  P_n^{|m|}(\cos\theta)\,e^{im\varphi},
\]
for \(n \geq 0\), \(-n \leq m \leq n\).  The spherical harmonics form a complete
orthonormal basis of \(L^2(S^2)\).
\end{definition}

%% ============================================================
\section{Hermite and Laguerre Polynomials}
\label{sec:hermite_laguerre}
%% ============================================================

\subsection{Hermite Polynomials}

\begin{definition}[Hermite polynomials, physicist's convention]
\[
  H_n(x) = (-1)^n e^{x^2}\frac{d^n}{dx^n}e^{-x^2}.
\]
\end{definition}

The first few are \(H_0 = 1\), \(H_1 = 2x\), \(H_2 = 4x^2-2\), \(H_3 = 8x^3-12x\).

\begin{theorem}[Orthogonality]
\[
  \int_{-\infty}^\infty H_n(x)H_m(x)e^{-x^2}\,dx = 2^n\,n!\,\sqrt{\pi}\,\delta_{nm}.
\]
\end{theorem}

\begin{theorem}[Recurrence relation]
\[
  H_{n+1}(x) = 2x\,H_n(x) - 2n\,H_{n-1}(x).
\]
\end{theorem}

\subsection{Quantum Harmonic Oscillator}

The stationary Schr\"odinger equation for the harmonic oscillator,
\[
  -\frac{\hbar^2}{2m}\psi'' + \frac{1}{2}m\omega^2 x^2\psi = E\psi,
\]
reduces via \(\xi = \sqrt{m\omega/\hbar}\,x\) to
\[
  \psi'' - (\xi^2 - \lambda)\psi = 0, \qquad \lambda = 2E/(\hbar\omega).
\]
Square-integrable solutions exist only for \(\lambda = 2n+1\), \(n \in \mathbb{N}_0\),
giving quantised energies \(E_n = \hbar\omega(n+\tfrac{1}{2})\) and wavefunctions
\[
  \psi_n(x) = \left(\frac{m\omega}{\pi\hbar}\right)^{1/4}
  \frac{1}{\sqrt{2^n n!}}\,H_n\!\left(\sqrt{\frac{m\omega}{\hbar}}\,x\right)
  e^{-m\omega x^2/(2\hbar)}.
\]

\subsection{Laguerre Polynomials}

\begin{definition}[Laguerre polynomials]
\[
  L_n(x) = \frac{e^x}{n!}\frac{d^n}{dx^n}(x^n e^{-x}) = \sum_{k=0}^n \binom{n}{k}
  \frac{(-x)^k}{k!}.
\]
\end{definition}

\begin{theorem}[Orthogonality]
\[
  \int_0^\infty L_n(x)L_m(x)e^{-x}\,dx = \delta_{nm}.
\]
\end{theorem}

\begin{theorem}[Recurrence]
\[
  (n+1)L_{n+1}(x) = (2n+1-x)L_n(x) - n\,L_{n-1}(x).
\]
\end{theorem}

\subsection{Hydrogen Atom}

The radial Schr\"odinger equation for the hydrogen-like atom after separation of variables
becomes, with appropriate substitutions, the associated Laguerre equation.  The radial
wavefunctions are
\[
  R_{nl}(r) = -\sqrt{\left(\frac{2}{na_0}\right)^3\frac{(n-l-1)!}{2n[(n+l)!]^3}}\,
  e^{-r/(na_0)}\!\left(\frac{2r}{na_0}\right)^l L_{n-l-1}^{2l+1}\!\!\left(\frac{2r}{na_0}\right),
\]
where \(a_0\) is the Bohr radius and \(L_n^\alpha\) are the generalised Laguerre polynomials
\(L_n^\alpha(x) = e^x x^{-\alpha}(d^n/dx^n)(e^{-x}x^{n+\alpha})/n!\).

%% ============================================================
\section{Elliptic Functions}
\label{sec:elliptic}
%% ============================================================

\subsection{Elliptic Integrals}

\begin{definition}[Complete elliptic integrals]
For modulus \(0 \leq k < 1\),
\begin{align*}
  K(k) &= \int_0^{\pi/2}\frac{d\theta}{\sqrt{1-k^2\sin^2\theta}}, \\
  E(k) &= \int_0^{\pi/2}\sqrt{1-k^2\sin^2\theta}\,d\theta.
\end{align*}
\end{definition}

\begin{theorem}[Legendre relation]
Let \(k' = \sqrt{1-k^2}\) be the complementary modulus.  Then
\[
  K(k)\,E(k') + K(k')\,E(k) - K(k)\,K(k') = \frac{\pi}{2}.
\]
\end{theorem}

\subsection{Weierstra\ss{} Elliptic Function}

\begin{definition}[Weierstra\ss{} \(\wp\)-function]
Let \(\Lambda = \omega_1\mathbb{Z} + \omega_2\mathbb{Z}\) be a lattice in \(\mathbb{C}\).
\[
  \wp(z; \Lambda) = \frac{1}{z^2}
  + \sum_{\omega \in \Lambda\setminus\{0\}} \left(\frac{1}{(z-\omega)^2} - \frac{1}{\omega^2}\right).
\]
\end{definition}

\begin{theorem}[Differential equation for \(\wp\)]
\[
  (\wp')^2 = 4\wp^3 - g_2\,\wp - g_3,
\]
where the Eisenstein series \(g_2 = 60 G_4(\Lambda)\) and \(g_3 = 140 G_6(\Lambda)\).
The discriminant \(\Delta = g_2^3 - 27g_3^2 \neq 0\) ensures the cubic has distinct roots.
\end{theorem}

\begin{remark}
The equation \(y^2 = 4x^3 - g_2 x - g_3\) is the Weierstra\ss{} form of an elliptic
curve.  Every elliptic curve over \(\mathbb{C}\) is isomorphic to \(\mathbb{C}/\Lambda\)
for some lattice \(\Lambda\), and the parametrisation \(z \mapsto (\wp(z), \wp'(z))\)
realises this isomorphism explicitly.
\end{remark}

\subsection{Jacobi Elliptic Functions}

The Jacobi functions \(\operatorname{sn}(u,k)\), \(\operatorname{cn}(u,k)\),
\(\operatorname{dn}(u,k)\) are defined as the inverses of the elliptic integral of the
first kind.  They satisfy
\begin{align*}
  \operatorname{sn}^2 + \operatorname{cn}^2 &= 1, \\
  k^2\operatorname{sn}^2 + \operatorname{dn}^2 &= 1, \\
  \frac{d}{du}\operatorname{sn}(u,k) &= \operatorname{cn}(u,k)\,\operatorname{dn}(u,k).
\end{align*}
The periods are \(4K(k)\) for \(\operatorname{sn}\) and \(\operatorname{cn}\), and
\(2K(k)\) for \(\operatorname{dn}\).

%% ============================================================
\section{Riemann Zeta and Dirichlet \texorpdfstring{$L$}{L}-Functions}
\label{sec:zeta}
%% ============================================================

\subsection{The Riemann Zeta Function}

\begin{definition}[Riemann zeta function]
For \(\Re(s) > 1\),
\[
  \zeta(s) = \sum_{n=1}^\infty \frac{1}{n^s} = \prod_p \frac{1}{1-p^{-s}},
\]
the product running over all rational primes \(p\).
\end{definition}

\subsection{Analytic Continuation and Functional Equation}

\begin{theorem}[Functional equation of \(\zeta\)]\label{thm:zeta_functional}
The function \(\zeta(s)\) extends to a meromorphic function on \(\mathbb{C}\) with a
single simple pole at \(s = 1\) with residue \(1\).  Defining the completed zeta function
\[
  \xi(s) = \tfrac{1}{2}s(s-1)\pi^{-s/2}\Gamma\!\left(\tfrac{s}{2}\right)\zeta(s),
\]
the functional equation states
\[
  \xi(s) = \xi(1-s).
\]
Equivalently,
\[
  \zeta(s) = 2^s\pi^{s-1}\sin\!\left(\frac{\pi s}{2}\right)\Gamma(1-s)\,\zeta(1-s).
\]
\end{theorem}

\begin{proof}[Proof sketch]
Starting from the Mellin transform \(\pi^{-s/2}\Gamma(s/2)\zeta(s)
= \int_0^\infty x^{s/2-1}\theta(x)\,dx\) where
\(\theta(x) = \sum_{n=1}^\infty e^{-\pi n^2 x}\), the modular transformation
\(\theta(1/x) = x^{1/2}\theta(x) + (x^{1/2}-1)/2\) (consequence of the Jacobi
theta-function identity) yields the symmetry \(s \leftrightarrow 1-s\) after splitting
the integral at \(x = 1\).
\end{proof}

\subsection{Special Values}

\begin{theorem}[Euler's formula for even arguments]
For \(n \in \mathbb{N}\),
\[
  \zeta(2n) = \frac{(-1)^{n+1}(2\pi)^{2n}B_{2n}}{2\,(2n)!},
\]
where \(B_{2n}\) are the Bernoulli numbers.  In particular,
\(\zeta(2) = \pi^2/6\), \(\zeta(4) = \pi^4/90\), \(\zeta(6) = \pi^6/945\).
\end{theorem}

\subsection{Dirichlet \(L\)-Functions}

\begin{definition}[Dirichlet character and \(L\)-function]
A Dirichlet character modulo \(q\) is a completely multiplicative function
\(\chi: \mathbb{Z} \to \mathbb{C}\) with \(\chi(n+q) = \chi(n)\) and \(\chi(n) = 0\)
if \(\gcd(n,q) > 1\).  The associated \(L\)-function is
\[
  L(s,\chi) = \sum_{n=1}^\infty \frac{\chi(n)}{n^s}
  = \prod_p \frac{1}{1-\chi(p)p^{-s}}, \qquad \Re(s) > 1.
\]
\end{definition}

\begin{theorem}[Non-vanishing at \(s=1\)]
For any non-principal character \(\chi\), \(L(1,\chi) \neq 0\).  This is the key
analytic input for Dirichlet's theorem on primes in arithmetic progressions.
\end{theorem}

\begin{theorem}[Functional equation for \(L(s,\chi)\)]
Let \(\chi\) be a primitive character modulo \(q\) with root number \(\varepsilon(\chi)\).
Define
\[
  \Lambda(s,\chi) = \left(\frac{q}{\pi}\right)^{(s+a)/2}\!\Gamma\!\left(\frac{s+a}{2}\right)L(s,\chi),
\]
where \(a = 0\) if \(\chi(-1) = 1\) and \(a = 1\) if \(\chi(-1) = -1\).  Then
\[
  \Lambda(s,\chi) = \frac{\tau(\chi)}{i^a\sqrt{q}}\,\Lambda(1-s,\bar\chi),
\]
where \(\tau(\chi) = \sum_{n=1}^q \chi(n)e^{2\pi in/q}\) is the Gauss sum.
\end{theorem}

%% ============================================================
\section{Applications}
\label{sec:applications}
%% ============================================================

\subsection{Quantum Mechanics}

Special functions permeate quantum mechanics.  The time-independent Schr\"odinger equation
in spherically symmetric potentials separates into a radial and an angular part:
\begin{itemize}
  \item \textbf{Angular part}: spherical harmonics \(Y_n^m(\theta,\varphi)\)
        (Legendre polynomials + exponentials).
  \item \textbf{Harmonic oscillator}: Hermite polynomials \(H_n\).
  \item \textbf{Hydrogen atom}: associated Laguerre polynomials \(L_{n-l-1}^{2l+1}\).
  \item \textbf{Linear potential / WKB}: Airy functions \(\operatorname{Ai}(x)\),
        \(\operatorname{Bi}(x)\).
  \item \textbf{Scattering}: Bessel and spherical Bessel functions \(j_l(kr)\).
\end{itemize}

\subsection{Classical Electrodynamics}

The Laplace equation \(\nabla^2\Phi = 0\) in cylindrical coordinates admits solutions in
terms of Bessel functions \(J_n(kr)\) and \(Y_n(kr)\).  In spherical coordinates the
solutions involve spherical harmonics.

\subsection{Number Theory and the Distribution of Primes}

The Prime Number Theorem states that \(\pi(x) \sim x/\ln x\) as \(x \to \infty\).
Its proof hinges on the non-vanishing of \(\zeta(s)\) on the line \(\Re(s) = 1\) and
the explicit formula
\[
  \psi(x) = x - \sum_\rho \frac{x^\rho}{\rho} - \ln(2\pi) - \tfrac{1}{2}\ln(1-x^{-2}),
\]
where the sum runs over the non-trivial zeros \(\rho\) of \(\zeta\).

\begin{conjecture}[Riemann Hypothesis]
All non-trivial zeros of \(\zeta(s)\) satisfy \(\Re(\rho) = 1/2\).
This conjecture remains one of the Clay Millennium Problems and is currently unproved.
\end{conjecture}

\subsection{Elliptic Functions and Number Theory}

The \(j\)-invariant
\[
  j(\tau) = 1728\,\frac{g_2^3}{\Delta} = q^{-1} + 744 + 196884\,q + \cdots,
  \quad q = e^{2\pi i\tau},
\]
is the hauptmodul for the modular group \(\mathrm{SL}(2,\mathbb{Z})\) and plays a central
role in the theory of complex multiplication and in Monstrous Moonshine.

%% ============================================================
\section{References}
%% ============================================================

\begin{thebibliography}{99}

\bibitem{AS}
M.~Abramowitz and I.~A.~Stegun,
\textit{Handbook of Mathematical Functions with Formulas, Graphs, and Mathematical Tables},
Dover Publications, New York, 1972.

\bibitem{NIST}
F.~W.~J.~Olver, D.~W.~Lozier, R.~F.~Boisvert, and C.~W.~Clark (eds.),
\textit{NIST Digital Library of Mathematical Functions},
Release 1.2.0, \url{https://dlmf.nist.gov/}, 2024.

\bibitem{WhittakerWatson}
E.~T.~Whittaker and G.~N.~Watson,
\textit{A Course of Modern Analysis}, 5th ed.,
Cambridge University Press, 2021 (reprint of 1927 edition).

\bibitem{Lebedev}
N.~N.~Lebedev,
\textit{Special Functions and Their Applications},
Dover Publications, New York, 1972.

\bibitem{Stein}
E.~M.~Stein and R.~Shakarchi,
\textit{Complex Analysis},
Princeton University Press, 2003.

\bibitem{GradshteynRyzhik}
I.~S.~Gradshteyn and I.~M.~Ryzhik,
\textit{Table of Integrals, Series, and Products}, 8th ed.,
Academic Press, 2015.

\bibitem{Andrews}
G.~E.~Andrews, R.~Askey, and R.~Roy,
\textit{Special Functions},
Cambridge University Press, 2000.

\bibitem{Watson}
G.~N.~Watson,
\textit{A Treatise on the Theory of Bessel Functions}, 2nd ed.,
Cambridge University Press, 1944.

\bibitem{Silverman}
J.~H.~Silverman,
\textit{The Arithmetic of Elliptic Curves}, 2nd ed.,
Springer, 2009.

\end{thebibliography}

\end{document}
